\documentclass[../Main]{subfiles}
\begin{document}
    \chapter{Desarrollo}
        Para el análisis de algunas característica de una memoria EPROM se plantea mostrar las formas del cálculos de los siguientes aspectos \cite{documentoTarea1}:
        \vspace{0.3cm}
        
        Para una memoria del tipo EPROM con el modelo 27C64 y utilizando la datasheet realizar los siguientes cálculos:
        \begin{enumerate}
            \item Basado en la tabla de pines o diagrama eléctrico del chip (DIP/SOIC) indicar la cantidad de bits del bus de direcciones.
            \item Basado en la tabla de pines o diagrama eléctrico del chip (DIP/SOIC) indicar la cantidad de bits del bus de datos.
            \item Mostrar el cálculo de la cantidad de celdas que posee la memoria.
            \item Mostrar el cálculo de la capacidad que posee la memoria.
            \item Calcular la dirección inicial de la memoria haciendo uso de una tabla donde se representen los diferentes pines del bus de direcciones y valor para la primera posición de memoria.
            \item Calcular la dirección final de la memoria haciendo uso de una tabla donde se representen los diferentes pines del bus de direcciones y valor para la posición final de memoria.
        \end{enumerate}    

        A continuación la respuesta de cada punto:

        \section{Bus de direcciones}
            Utilizando como base el datasheet del fabricante, se pudo determinar mediante la Fig. \ref{Fig: Diagrama de bloques de la memoria} que la memoria 27C64 posee como terminales de entrada las etiquetas desde $A_0-A_{12}$.

            \begin{figure}[ht]
                \centering
                \includegraphics[width=0.6\linewidth]{Figuras/Diagrama de bloques.png}
                \caption{Diagrama de bloques de la memoria}
                \label{Fig: Diagrama de bloques de la memoria}
            \end{figure}

            Según este diagrama y el análisis de los pines, puede concluirse que la memoria posee 13 bits de bus de direcciones.

        \section{Bus de datos}
            Sobre el concepto de bus de datos, recordar que estas terminales se pueden reconocer con la palabra O, del inglés \textit{output}, de manera que mediante la Fig. \ref{Fig: Diagrama de bloques de la memoria} se puede reconocer, que la memoria en análisis posee 8 bits de bus de datos.

        \section{Cantidad de celdas}
            Para lograr conocer la cantidad de celdas de una arquitectura puede utilizarse como formula general la Ec. \ref{Ec: Ecuacion sobre combinaciones binarias}, la cual posee como base que n es la cantidad de bits del bus de direcciones.

            \begin{equation}
                \text{Cantidad de celdas} = 2^n
                \label{Ec: Ecuacion sobre combinaciones binarias}
            \end{equation}

            Para el ejemplo en análisis, el bus de direcciones es de 13 bits, por tanto el cálculo es:
            \begin{equation}
                \text{Cantidad de celdas} = 2^{13}=8192
            \end{equation}   

            Como resultado, se puede interpretar que la memoria EPROM 27C64 posee 8192 celdas, representado en una matriz con un total de 512 filas con 16 columnas.

        \section{Capacidad de la memoria}
            Esta variable demanda el uso del concepto de bus de direcciones y de bus de datos, dado que es importante conocer la cantidad de celdas de la memoria y el tamaño del dato que posee cada celda, para realizar este calculo se utiliza la Ec. \ref{Ec: Ecuación de la capacidad}.

            \begin{equation}
                \text{Capacidad de memoria}=\text{Celdas totales x Ancho del bus en bits}
                \label{Ec: Ecuación de la capacidad}
            \end{equation}

            Con el Ec. \ref{Ec: Ecuación de la capacidad} puede realizarse el calculo con los datos antes recopilados.

            \begin{equation}
                Capacidad=2^n \cdot m
            \end{equation}

            \begin{equation}
                Capacidad=2^{13} \cdot \text{8 bits}
            \end{equation}

            \begin{equation}
                Capacidad=8KB
            \end{equation}

            La memoria 27C64 posee una capacidad de 8KB, esto como resultado de analizar la cantidad de bits del bus direcciones y del bus de datos.
            
        \section{Direcciones}
            Para conocer las direcciones principales de una memoria, la primera y la ultima permiten generar la idea global, por este motivo en la Figura \ref{Fig: Datos de la primera y ultima dirección} se muestra como se obtuvo la primera dirección de memoria, la cual se calcula conociendo que todos los pines del bus de direcciones de esa memoria se colocan con valor 0b y para la última dirección se establecen valores 1b.

            \begin{figure}[H]
                \centering
                \includegraphics[width=1\linewidth]{Figuras/Tabla.png}
                \caption{Datos de la primera y ultima dirección}
                \label{Fig: Datos de la primera y ultima dirección}
            \end{figure}

            Los resultados obtenidos son:
            \begin{itemize}
                \item Primea dirección de memoria es 0000H.
                \item Ultima dirección de memoria es 1FFFH
            \end{itemize}

        Con todos los cálculos mostrados, es posible tener una idea global de como esta compuesta la memoria 27C64.
\end{document}
